\section{Conclusion}
\label{sec:conclusion}

This section first summarises what the model comparisons establish, then states what the mechanism tests leave unresolved, and finally identifies the next steps.

This study built an edge-level intrusion-detection pipeline with structural and semantic representations of the same aggregated traffic. It compared a GNN model, a semantic prototype model, a fusion model and a feedback model on NF-UNSW-NB15 and NF-ToN-IoT. The feedback model was evaluated with two consultants, and the trained semantic head was also measured alone. Evaluation used pooled five-fold out-of-fold predictions at three training seeds and a fixed fold partition.

The trained-head feedback configuration is separated above the GNN, semantic and fusion stages and the faithful re-trained baselines on both datasets. However, the trained head alone has slightly higher mean macro-F1 on both, and its difference from feedback is not separated. The results support using the trained head as a stronger semantic consultant than the prototype readout. They do not establish an additional benefit from the complete feedback design over that consultant. The consultant change also changed the consultation percentile, so the increase between configurations cannot be assigned to either change alone.

The mechanism diagnostics reinforce the need for this distinction. With output fusion disabled, neither realistic consultant provides a separated injection gain at the earlier consultation settings tested, while an oracle with access to true labels does. Attention injection changes no predicted classes between passes in the tested configurations. Fusion alone is not separated from the GNN on either dataset. These negative results are part of the contribution: adding another interaction between the branches does not by itself establish a better classifier.

The conclusions remain within the controlled graph-construction scope defined in Section~\ref{sec:method-graph}. The next steps are matched mechanism tests at the final settings, broader uncertainty measurement, controls for message passing and the weighted branch split, and an evaluation built from label-independent traffic aggregation. The current evidence establishes the value of the trained semantic classifier and clearly identifies what remains unproven about feedback.
